\documentclass[11pt]{article}
\usepackage{preamble}
\setlength{\parskip}{4pt}

\begin{document}

\daybanner{AP Statistics \textperiodcentered\ Unit 1 \textperiodcentered\ Topic 1.11 \textperiodcentered\ B: 2026-09-29 \textperiodcentered\ E: 2026-10-07 \textperiodcentered\ TEACHER KEY}{Three Ways to Sample Students}

\sectionbanner{CED TETHER}

{\small
\begin{itemize}[leftmargin=1.5em,itemsep=2pt,topsep=2pt]
  \item \textbf{Skill 1.C} --- Describe an appropriate method for gathering and representing data.
  \item \textbf{EU DAT-2} --- The way we collect data influences what we can and cannot say about a population.
  \item \textbf{LO DAT-2.C} --- Identify a sampling method, given a description of a study. [Skill 1.C]
  \item \textbf{EK DAT-2.C.1} --- When an item from a population can be selected only once, this is called sampling without replacement. When an item from the population can be selected more than once, this is called sampling with replacement.
  \item \textbf{EK DAT-2.C.2} --- A simple random sample (SRS) is a sample in which every group of a given size has an equal chance of being chosen. Mechanisms include numbering individuals and using a random number generator, a table of random numbers, or drawing a card from a deck without replacement.
  \item \textbf{EK DAT-2.C.3} --- A stratified random sample involves dividing a population into separate groups called strata based on shared attributes (homogeneous grouping). Within each stratum a simple random sample is selected, and the selected units are combined to form the sample.
  \item \textbf{EK DAT-2.C.4} --- A cluster sample involves dividing a population into smaller groups called clusters, ideally with heterogeneity within each cluster and similarity among clusters. A simple random sample of clusters is selected, and data are collected from all observations in the selected clusters.
  \item \textbf{EK DAT-2.C.5} --- A systematic random sample selects members from a population according to a random starting point and a fixed, periodic interval.
  \item \textbf{EK DAT-2.C.6} --- A census selects all items/subjects in a population.
  \item \textbf{LO DAT-2.D} --- Explain why a particular sampling method is or is not appropriate for a given situation. [Skill 1.C]
  \item \textbf{EK DAT-2.D.1} --- There are advantages and disadvantages for each sampling method depending upon the question to be answered and the population from which the sample will be drawn.
\end{itemize}
}
\textbf{Essential question:} Which random sampling method best matches the structure of the population and the question?\par
\textbf{DOK-3 skill:} 4.A \textperiodcentered\ \textbf{FRQ pattern:} {\small\texttt{srs-\allowbreak{}vs-\allowbreak{}systematic-\allowbreak{}vs-\allowbreak{}stratified-\allowbreak{}choose-\allowbreak{}and-\allowbreak{}defend}}\par
\textbf{DOK rationale:} Part (c) selects a sampling method by linking grade differences in sleep to representation, then explains why randomness does not make the three plans equivalent.\par

\frameworkphaseheader{Do Now (first take) $\rightarrow$ Explore (video + follow-along) $\rightarrow$ Exit (finish + turn in)}{1 $\rightarrow$ 3}{45}{%
\begin{itemize}[leftmargin=1.5em,itemsep=2pt,topsep=2pt]
  \item Board slide up at the bell. Ask students to compare any two methods before naming them.
  \item Begin the video follow-along at minute 5.
  \item Return to the board slide at minute 33. Circulate and ask what each method guarantees.
  \item Collect at the door. Score part (c) only, E/P/I.
\end{itemize}
}{%
\begin{itemize}[leftmargin=1.5em,itemsep=2pt,topsep=2pt]
  \item Commit to whether the methods are equivalent and cite one difference.
  \item Complete the video follow-along about SRS, systematic, stratified, and cluster sampling.
  \item Finish (a)--(c), revise the first take in the exit line, and turn in the sheet.
\end{itemize}
}{%
\begin{itemize}[leftmargin=1.5em,itemsep=2pt,topsep=2pt]
  \item What does Method 3 require in every sample?
  \item Could the hat produce no students from one grade?
  \item Does a random start make every group of 50 possible?
  \item Why does grade matter for this question?
\end{itemize}
}{Solo teacher. Circulate ~5 pairs. Ask for the guarantee and the failure mode before accepting a method name; do not treat all chance-based designs as interchangeable.}

\begin{callout}{\IconWarn\ WATCH FOR}{calloutred}
\begin{itemize}[leftmargin=1.5em,itemsep=2pt,topsep=2pt]
  \item Calling every random method an SRS
  \item Calling sampling within every grade cluster sampling
  \item Claiming an SRS guarantees equal grade counts
\end{itemize}
\end{callout}

\sectionbanner{SCORING PART (c) --- E / P / I}

\textbf{Expected elements}
\begin{itemize}[leftmargin=1.5em,itemsep=2pt,topsep=2pt]
  \item \textbf{grade-based-choice} (required) --- Chooses Method 3 because grade is related to sleep
  \item \textbf{guarantee-and-contrast} (required) --- Explains 12 or 13 per grade versus no required grade counts in Methods 1 and 2
  \item \textbf{rejects-equivalence} (required) --- Explains that use of chance does not make the possible samples or guarantees identical
\end{itemize}

{\small\renewcommand{\arraystretch}{1.25}
\noindent\begin{tabular}{|p{0.5in}|p{5.9in}|}
\hline
\textbf{E} & All three elements are correct and linked to this problem. Accept equivalent plain-language reasoning. \\ \hline
\textbf{P} & At least one element includes a correct evidence-to-reason link, but another is missing or incorrect. \\ \hline
\textbf{I} & No correct evidence-to-reason link; only a choice, copied numbers, or unsupported statements. \\ \hline
\end{tabular}}

\textbf{Common mistakes}
\begin{itemize}[leftmargin=1.5em,itemsep=2pt,topsep=2pt]
  \item Calling every random method an SRS
  \item Calling sampling within every grade cluster sampling
  \item Claiming an SRS guarantees equal grade counts
\end{itemize}

\clearpage
\focusbanner{TODAY'S PROBLEM --- annotated}

Suppose a school with four grades of 300 students each wants to estimate nightly hours of sleep. Three teams propose samples of 50.\par\smallskip \textbf{Method 1:} Choose a random start among the first 24 names on the alphabetical roster, then take every 24th name.\par \textbf{Method 2:} Put all 1{,}200 names in a hat, mix thoroughly, and draw 50 without replacement.\par \textbf{Method 3:} Separate the roster by grade, then randomly choose 12 students from each grade, then choose two grades at random and add one more randomly chosen student from each, for 50 total.\par\smallskip A student says, \emph{They all use randomness, so they are the same.}\par
\begin{center}
\textbf{Three proposed samples}\par\smallskip
\begin{tabular}{|l|c|c|}
\hline
\textbf{\#} & \textbf{Selection rule} & \textbf{Total} \\ \hline
\textbf{1} & Every 24th & 50 \\ \hline
\textbf{2} & Names from hat & 50 \\ \hline
\textbf{3} & Names within grades & 50 \\ \hline
\end{tabular}
\end{center}
\begin{firsttakebox}{FIRST TAKE (5 min)}
Would these plans give the school the same mix of students every time? Give one reason from the plans.\par
\answer{Not graded. Expect the word random to make the three methods look equivalent at first.}
\end{firsttakebox}

\textit{Rules callout ``THREE RANDOM METHODS (keep on the page)'' is printed on the student sheet.}\par\medskip

\dokbadge{1}~\textbf{(a)} Name each method: systematic random, simple random, or stratified random.\par
\answer{Method 1 is systematic random sampling. Method 2 is a simple random sample. Method 3 is stratified random sampling, with grade as the stratum.}

\dokbadge{2}~\textbf{(b)} Which method is an SRS of 50 students? Explain using what must have an equal chance, not just the word random.\par
\answer{Method 2 is the SRS because every possible group of 50 students has an equal chance to be drawn. Method 1 allows only groups generated by one of 24 starts, and Method 3 forces 12 or 13 students from each grade, so neither gives every group of 50 an equal chance.}

\dokbadge{3}~\textbf{(c)} Suppose sleep differs a lot by grade. Choose a method and explain what it guarantees about the mix of grades. Contrast it with the other two methods, then respond to the student's claim that all three are the same.\par
\answer{Choose Method 3: it guarantees 12 or 13 randomly selected students from each grade, so a grade with different sleep habits cannot be left out. Methods 1 and 2 do not fix the number from each grade, even though both use chance. All three can be random without allowing the same possible groups or guaranteeing the same grade mix.}

\begin{summaryexitbox}{EXIT}
One thing the video changed about my first take: \hrulefill
\end{summaryexitbox}

\end{document}
